\documentclass[a4paper,12pt]{ltjsarticle}

% パッケージ
\usepackage{graphicx}
\usepackage{amsmath, amssymb}
\usepackage{xcolor}
\usepackage{geometry}
\usepackage{enumitem}
\usepackage{hyperref}
\usepackage{luatexja-fontspec}

% フォント設定
\setmainjfont{IPAexMincho}
\setmainfont{Times New Roman}
\setsansjfont{IPAexMincho}
\setsansfont{Times New Roman}

% レイアウト
\geometry{margin=25mm}
\setlist[itemize]{label=--, leftmargin=2em}
\setlist[enumerate]{leftmargin=2em}

%====================
% 強調用コマンド
%====================
\newcommand{\important}[1]{\textcolor{blue}{\textbf{#1}}}

\begin{document}

%====================
% タイトルブロック
%====================
\begin{center}
  {\huge 進捗報告（ver.37）}\\[1em]
  {\large \today}\\
  {\large 千葉工業大学大学院 学習工学研究室}\\
  {\large 内山 裕太}
\end{center}

\vspace{2em}

%====================
\section{進捗概要}
\begin{itemize}
  \item 数学学習の分野調べ
  \item 自身の研究テーマの確定
  \item それっぽいシステムの開発（これから変えていく）
\end{itemize}
\newpage

%====================
\section{進捗詳細}
\subsection{数学学習の分野調べ}
数学学習を題材とした研究がどのくらい，どのようにされているのかを調査した．結果としては，\href{https://www.jstage.jst.go.jp/browse/jasme/29/0/_contents/-char/ja}{\important{全国数学教育学会}}に面白そうな研究がたくさんあることが分かった．\\

この学会では「研究」をメインに扱っているらしく，別の組織で\important{日本数学教育学会}もあるらしい．こちらは実践も重要視しているらしく，実際の数学教員なども多く在籍しているらしい．\\

今回は全国数学教育学会の学会誌（数学教育学研究），日本数学教育学会誌の中から以下のものをリストアップした

\begin{enumerate}
  \item \href{https://www.jstage.jst.go.jp/article/jjsme/65/3/65_13/_pdf/-char/ja}{中西知真紀, 国宗進, 家田晴行, 榎戸章仁, 春日龍郎, 金児正史, 小関熙純, 東風谷幸広, 山下国広: 図形における論証指導について, 日本数学教育学会誌, Vol. 65, No. 3, pp. 13-24, 1983}
  \item \href{https://www.jstage.jst.go.jp/article/jasme/27/1/27_33/_pdf/-char/ja}{渡邊慶子, 岡崎正和: 証明言語の生成とふり返りによる定理と証明の相互理解 ―場合分けのある証明に着目して―, Vol. 27, No. 1, pp.33-46, 2021}
  \item \href{https://www.jstage.jst.go.jp/article/jjsme/104/1/104_2/_pdf/-char/ja}{近藤裕: 算数・数学科における「説明・証明」の能力に関する研究, Vol. 104, No. 1, pp. 2-12, 2022}
  \item \href{https://www.jstage.jst.go.jp/article/jasme/3/0/3_KJ00008199223/_pdf/-char/ja}{長谷川順一, 三井誠弦: 中学生の証明問題解決過程の分析, Vol. 3, pp. 137-146, 1997}
  \item \href{https://www.jstage.jst.go.jp/article/jjet/42/Suppl./42_S42092/_pdf/-char/ja}{濱田さとみ, 鷹岡亮, 横山誠: 中学校数学科合同証明を対象とした証明構造理解支援Webアプリの開発と有用性検討, 日本教育工学会論文誌, Vol. 42, No. Suppl, pp. 169-172, 2018}
  \item \href{https://www.jstage.jst.go.jp/article/jjsme/69/11/69_23/_pdf/-char/en}{磯田正美: 体系化の立場から見た中2の図形指導, 日本数学教育学会誌, Vol. 69, No. 11, pp. 23-32}
  \item \href{https://www.jstage.jst.go.jp/article/jasme/23/2/23_141/_pdf/-char/ja}{宍戸建太, 岡崎正和: 図形の証明の構成過程におけるジェスチャーの役割に関する研究, 全国数学教育学会誌 数学教育学研究, Vol. 23, No. 2, pp. 141-149, 2017}
  \item \href{https://www.jstage.jst.go.jp/article/jjsme/77/11/77_17/_pdf/-char/ja}{垣花京子, 清水克彦: 図形の証明問題での測定値の役割 -コンピュータ環境下における生徒の活動分析を通して-, 日本数学教育学会誌, Vol. 77, No. 11, pp. 17-22}
\end{enumerate}

現在は上記の文献の選定（中身まで確認できているわけではない）だけであり，軽く目を通すこともできていない状態．ただ，自分のやりたいことと今までの経験からうまく結びつけられそうな分野を見つけた．

\subsection{自身の研究テーマの確定}
今まで「知識の構造化」という考え方で研究をやってきたが，部品的知識の考え方と構造的な考え方を使いこなす必要のある分野を見つけた．\\

中学校2年生から始まる「図形の証明学習」\\

図形の証明学習では，学習者の持つ知識（図形の定義や性質）を「既知部品」として，問題から与えられる仮定などを「一時部品」とすることができそう．また，結論を導くために必要な3つの部品が何かを考え，それらが「自分の持ち合わせる部品のみで構築できるのか」や「部品同士を組み合わせて新たな部品を生み出す必要があるのか」を考える必要が出てくる．意外と，定義や性質，合同条件などの暗記情報は覚えてくれる子が多い一方で，例題ベースの証明問題以外は解けないとあきらめる子が多い印象がある．\\

上記の考え方ベースで学べる教材があれば，証明問題は「組み立てるだけ」という考えで学ぶことができるはずなので，有用性はあると考えました（中学生は「論理的に」証明をするのは難しすぎるので，「結論には何が必要か」を考える力が重要になると思います）．\\

という考え方の下でJSiSEの学会原稿を書いてみようと思ったですが，まだ手法が固まってないのと，関連研究が読み切れてないから難しかったです．自分なりに研究背景や問題意識を整理して1章だけは書いてみました．一緒に添付しておきます．\\

この研究テーマに対して，少し意見がいただきたいです．僕は結構良い案だと思っているのですが，有用性が本当にあるのかは皆さんにお聞きしたいです．
\newpage

\subsection{それっぽいシステムの開発}
2.2の内容を基に，自分の想像するシステムを簡単に構成してみました．ChatGPTに相談し，TypeScript（JavaScript）の勉強もかねてそれっぽいのができました．\\

\includegraphics[width=\linewidth]{figs/20260306/system1.png}

今回の仮開発を通して，画面の構成案はなんとなく見えました．あとは，どんな機能をどのように実装するかだと思ってます．基本的には，定義や性質などの「暗記情報」を学ぶフェーズも用意したいので，最初のフェーズで「部品獲得」をして，次のフェーズで「部品利用」という流れになりそうです．もちろん，問題によっては仮定等から部品を一時的に作り出すことがあると思います（持っている知識だけで解ける問題だけではないですから）．少し話に上げていた「問題文から必要な情報を読み取る能力」についても，考えていく必要がありそうです．
\newpage

%====================
\section{問題点・課題}
\begin{itemize}
  \item 日々の生活に対するモチベが終わってた（1週間ほぼ寝込んでた...）
  \item 全国大会の原稿を4月中に一旦書きたい（手法の提案になるかなぁ...）
\end{itemize}

%====================
\section{今後の予定}
\begin{enumerate}
  \item システムの機能要件作成
  \item タスクの細分化
  \begin{itemize}
    \item システム機能要件の要件別のタスク化
    \item 研究室での活動を見える化するための仕組みづくり（済）
    \item 研究に飽きたときの勉強決め
  \end{itemize}
\end{enumerate}
\newpage

%====================
\section*{雑談}
なぜあのレベルまで気力がなくなってしまったのだろうか...と思うくらい何もする気力がわかなかったです...\\

多分，「管理されないとうまく動けない」と「目標が見えていないとうまく動けない」という2つの性質を持っているので，言葉で表すのが難しいですが，何かが爆発してしまったんだと思います．自分なりに改善しようと思い，日々の活動やタスク内容を管理できる枠組みを作りました．また，毎日決まった時間にスマホ（Macも）に通知が来るように設定して，スケジュール通りの活動ができるような環境を作りました．具体的には，タイムカード的な仕組みを作って，その日の活動時間や行ったタスクをデータとしてまとめられるようにして，目に見える進捗を自分で褒められるようなものです．\\

また，最近勉強に対するモチベーションが上がってきているので，研究に飽きたときとかに色々な科目を勉強したいと思いました．

\begin{itemize}
  \item 古典
  \item 英語（本当に英語力がゴミです）
  \item 化学（エンタルピーは勉強した）
  \item 生物（はたらく細胞は神作です）
  \item 物理（理科の中で一番簡単ですよね？）
\end{itemize}

上記の内容を勉強するために，同じ職場の人とかに調査をしておすすめの参考書をリスト化しました．生徒用の参考書とかが校舎にあるので，それらを借りながら勉強してみます．\\

受験生を全員合格させたということは，僕は本当にラスト1年を切ったことになります．今年は9月までに研究をある程度のところまでもっていく！9月までに実験する！という意気込みで頑張ります！\\\\

ChatGPT-Plusに積み立てNISAの相談をしたら，すごい熱心に色々教えてくれました．10年後の僕に幸あれ！
\newpage

%====================
\section*{備考・メモ}
\begin{itemize}
  \item 3/26～4/3は春期講習（おそらく）\\
  \item 4月からの院ゼミは何曜日になりそうか\\（早めにわかれば，院ゼミの日はシフト抜いてくれるらしい）\\
  \item 全国大会出れるなら出たい（やってやります）
\end{itemize}

\end{document}
